\section{宏观 Scaling：从模型曲线到资本与价格信号}

AI 扩张常被一张曲线概括，但产业系统至少同时产生四类信号：模型能力随训练资源变化的技术信号，工具使用增长的产品信号，数据中心与芯片投资的资本信号，以及 GPU 租赁价格的供需信号。它们可以相互支持，却不是同一件事。技术进步不保证商业收入，收入增长不保证资本配置有效，资本投入也不保证每一块硬件都能被高效使用。

\subsection{模型能力的可预测性}

Scaling law 的工程价值在于预算和风险管理：团队可以用历史实验估计增加数据、参数或计算后损失会怎样变化，从而决定下一次训练规模。预测区间越远，架构变化、数据质量、后训练和评测漂移带来的误差越大。宏观图表适合展示长期方向，不能替代对单个模型、任务和训练 recipe 的审计。

\begin{figure}[H]
\centering
\includegraphics[width=0.94\textwidth]{images/00-44-00-macro-ai-scaling.jpg}
\caption{讲者用宏观模型扩张图说明，AI scaling 在若干历史区间内表现出可观察的规律。}
\figsource{Stanford Online 官方视频 00:44:00，`images/00-44-00-macro-ai-scaling.jpg`。}
\end{figure}

\begin{knowledgebox}{读图：区分观测点、拟合线与未来假设}
观测点来自已经公开或估算的训练项目，拟合线总结过去趋势，向右延伸的部分则包含未来假设。读者应先看纵横轴单位和数据来源，再问模型代际是否可比、后训练是否被计入、评测是否变化。拟合很好只能证明所选数据在该坐标系中规律明显，不能证明下一代系统一定沿同一斜率增长。
\end{knowledgebox}

\subsection{产品采用：Claude Code 的增长说明什么}

产品信号比 benchmark 更接近真实需求。讲者展示 Claude Code 在十二个月内的日提交增长，用来说明 coding agent 正迅速进入开发流程。这个指标至少意味着用户愿意把模型生成或协助的代码提交到仓库，但它仍不能单独回答代码质量、返工成本、安全事件、付费留存或不同团队的净生产率。

\begin{figure}[H]
\centering
\includegraphics[width=0.94\textwidth]{images/00-46-02-claude-code-growth.jpg}
\caption{Claude Code 日提交量在十二个月内快速增长，展示 agent 产品进入真实工作流的速度。}
\figsource{Stanford Online 官方视频 00:46:02，`images/00-46-02-claude-code-growth.jpg`。}
\end{figure}

\begin{importantbox}{读图：采用量不是生产率结论}
提交量可以证明使用强度上升，不能直接证明每次提交都提高净产出。完整评估还需要 merge rate、review time、缺陷率、回滚率、开发者覆盖和长期维护成本。产品增长仍然重要，因为它扩大了 context 与反馈来源；只是这些反馈必须经过质量与权限筛选，才能进入下一轮模型或工具改进。
\end{importantbox}

\subsection{Great Transition 的资本开支}

当应用需求、模型规模和算力预期一起上升，云厂商与基础设施公司会提前采购土地、电力、网络和加速器。\term{Capex} 是 capital expenditure，即用于长期资产的资本开支；它与日常运营费用不同，往往具有建设周期长、回收慢和退出困难的特点。AI 需求预测一旦偏离，过度建设与供应不足都可能持续多年。

\begin{figure}[H]
\centering
\includegraphics[width=0.94\textwidth]{images/00-46-45-great-transition.jpg}
\caption{讲者把近期基础设施资本开支上升描述为 Great Transition 的宏观信号。}
\figsource{Stanford Online 官方视频 00:46:45，`images/00-46-45-great-transition.jpg`。}
\end{figure}

\begin{knowledgebox}{读图：capex 是下注，不是已经实现的价值}
资本开支曲线表明企业相信未来会有足够的软件收入和战略价值支撑硬件投入。数据中心建成以后，价值还取决于利用率、模型效率、电力价格、互连能力和客户需求。把 capex 当作需求信心的量化表达是合理的；把它当作生产率或社会收益已经兑现的证明，则跨越了尚未验证的链条。
\end{knowledgebox}

\subsection{GPU 租赁价格：异构算力市场的温度计}

价格把供给、需求、地域、电力、硬件代际和合约期限压成一个数字，因此很有吸引力。讲者展示 AMP 观察到的 GPU 租赁价格变化，用它说明稀缺性并非单调上升：新供应进入、旧卡折价、软件效率改善和需求结构变化都会改变价格。更重要的是，不同 GPU 并不是同一种商品，标题价格必须和显存、互连、可靠性、可用时长及集群规模一起阅读。

\begin{figure}[H]
\centering
\includegraphics[width=0.94\textwidth]{images/00-47-45-gpu-rental-prices.jpg}
\caption{GPU 租赁价格随硬件代际与市场阶段变化，反映供需、折旧和可替代性的共同作用。}
\figsource{Stanford Online 官方视频 00:47:45，`images/00-47-45-gpu-rental-prices.jpg`。}
\end{figure}

\begin{warningbox}{读图：每 GPU-hour 仍不是统一商品价格}
同样标成一小时，单卡、八卡节点和大规模无阻塞集群的价值完全不同；短租、预留、抢占与长期合同也不可直接比较。训练还依赖网络拓扑、存储吞吐、软件栈与故障恢复。价格图能显示趋势，却不能在缺少规格与服务条件时证明“算力已经便宜”或“某卡一定更划算”。
\end{warningbox}

\subsection{本章小结}

模型曲线、产品采用、资本开支和租赁价格分别回答技术、需求、投资和供需问题。把它们并置可以观察转型是否跨层发生，但不能把四类信号合成一条必然因果链。下一节用历史基础设施周期寻找可复用模式，同时明确类比的边界。

